% ============================================================================
% Fichier  : partie4/P04_C14_forensique_mobile.tex : FORN (LPI342)
% Sources : NIST SP 800-101r1 ; ISO/IEC 27037. Opération TSANAGA.
% ============================================================================

\chapter{Forensique Mobile}
\label{chap:mobile}
\chaptermark{Forensique mobile}

\epigraph{%
    « Le téléphone est le témoin qui ne quitte jamais son propriétaire. »%
}{Principe du cours}

\niveauChapitre{Intermédiaire à expert}{Chapitre~\ref{chap:disque}~; notions
de systèmes mobiles.}

\begin{boxobjectifs}
\begin{itemize}
    \item Comprendre les spécificités de l'enquête sur smartphone.
    \item Identifier les traces (messages, applications, géolocalisation).
    \item Connaître les niveaux d'extraction.
    \item Respecter le cadre juridique du mobile.
    \item Qualifier un artefact mobile.
\end{itemize}
\end{boxobjectifs}

\bigskip

% ============================================================================
\section*{Question directrice : le téléphone de M. NDONGO confirme-t-il la revente des données ?}
% ============================================================================

Le téléphone saisi peut contenir messages, contacts, et traces de
négociation. L'enquête mobile obéit à des règles propres. Ce chapitre les
présente.

\medskip
\textbf{Plan du chapitre.}
\begin{enumerate}
    \item Les spécificités du mobile (section~\ref{sec:c14-spec}).
    \item Les traces mobiles (section~\ref{sec:c14-traces}).
    \item Les niveaux d'extraction (section~\ref{sec:c14-extraction}).
    \item Le cadre juridique (section~\ref{sec:c14-droit}).
\end{enumerate}

\separateur

% ============================================================================
\section{Les Spécificités du Mobile}
\label{sec:c14-spec}
% ============================================================================

Un smartphone diffère d'un ordinateur~: stockage chiffré par défaut,
verrouillage, connexion permanente (risque d'effacement à distance),
diversité des systèmes (Android, iOS). On l'isole des réseaux dès la saisie
(mode avion, sachet de blindage) pour éviter l'effacement.

% ============================================================================
\section{Les Traces Mobiles}
\label{sec:c14-traces}
% ============================================================================

\begin{tableaucours}{L{3.6cm}L{8.6cm}}{%
    \tableauHeader{Trace} &
    \tableauHeader{Ce qu'elle révèle}
}
Messages & SMS, messageries, conversations \\
Journaux d'appels & Qui, quand, combien de temps \\
Applications & Données locales, historiques, pièces jointes \\
Géolocalisation & Lieux fréquentés, déplacements \\
Métadonnées de photos & Date, lieu de prise de vue \\
\end{tableaucours}
\vspace{-0.8em}
\begin{center}{\small\itshape Tableau~14.1. Traces d'un smartphone~\cite{nist800101}}\end{center}

% ============================================================================
\section{Les Niveaux d'Extraction}
\label{sec:c14-extraction}
% ============================================================================

\begin{tableaucours}{L{3.6cm}L{8.6cm}}{%
    \tableauHeader{Niveau} &
    \tableauHeader{Portée}
}
Manuel & Lecture à l'écran, photographiée~; limité, mais simple \\
Logique & Données accessibles par l'interface (messages, contacts) \\
Physique & Copie bas niveau~; plus complet, plus complexe \\
\end{tableaucours}
\vspace{-0.8em}
\begin{center}{\small\itshape Tableau~14.2. Niveaux d'extraction mobile}\end{center}

% ============================================================================
\section{Le Cadre Juridique}
\label{sec:c14-droit}
% ============================================================================

\begin{boxCRO}
\textbf{Le mobile, intime par nature.} Un téléphone contient la vie privée
entière de son propriétaire~: messages intimes, santé, famille. L'examiner
sans limite dégrade la \cro{-}{0}{0}. On s'en tient strictement au périmètre
de la réquisition et l'on protège les tiers (correspondants).
\end{boxCRO}

\begin{boxscenario}{Opération TSANAGA. Les messages de négociation}
\textbf{Constatation.} L'extraction logique du téléphone, dans le périmètre,
révèle des messages datés de janvier évoquant la remise d'un « fichier » contre
paiement, à un contact extérieur.

\medskip
\textbf{Analyse} (confiance élevée pour les messages, moyenne pour leur
portée exacte). Les messages corroborent fortement la revente~; leur
interprétation précise (quel fichier, quel montant) se recoupe avec le disque
et le réseau. On rapporte les messages, à charge et à décharge.
\end{boxscenario}

\begin{boxmetier}{Enquêteur en cybercriminalité et expert judiciaire}
\textbf{Isoler, puis extraire dans le périmètre.} L'enquêteur isole le
téléphone des réseaux dès la saisie, choisit le niveau d'extraction adapté, et
limite l'examen au périmètre. Le respect de l'intimité conditionne la
recevabilité.
\end{boxmetier}

\separateur

\begin{boxresume}
\textbf{Ce qu'il faut retenir de ce chapitre~:}
\begin{itemize}
    \item \textbf{Le mobile} est chiffré, verrouillé, connecté~: on l'isole.
    \item \textbf{Les traces}~: messages, appels, applications,
          géolocalisation.
    \item \textbf{Trois niveaux d'extraction}~: manuel, logique, physique.
    \item \textbf{Le mobile est intime}~: on s'en tient au périmètre.
    \item \textbf{On recoupe} avec le disque et le réseau.
\end{itemize}
\end{boxresume}

\begin{boxinfo}
\textbf{Points de contrôle de l'élève.} Je sais~: expliquer les spécificités
du mobile~; citer les traces~; distinguer les niveaux d'extraction~; respecter
le périmètre~; qualifier un artefact mobile.
\end{boxinfo}

\bigskip

\begin{boxexo}[title={\ding{45}\quad Exercice~14.1. Spécificités \niveaubase}]
Pourquoi isole-t-on un téléphone des réseaux dès la saisie~?
\end{boxexo}

\begin{boxexo}[title={\ding{45}\quad Exercice~14.2. Traces \niveaubase}]
Citez quatre traces d'un smartphone.
\end{boxexo}

\begin{boxexo}[title={\ding{45}\quad Exercice~14.3. Niveaux \niveaubase}]
Distinguez extraction logique et physique.
\end{boxexo}

\begin{boxexo}[title={\ding{45}\quad Exercice~14.4. Effacement \niveaumoyen}]
Quel risque court un téléphone resté connecté après la saisie~?
\end{boxexo}

\begin{boxexo}[title={\ding{45}\quad Exercice~14.5. Intimité \niveaumoyen}]
Pourquoi le mobile exige-t-il une attention particulière à l'intimité~?
\end{boxexo}

\begin{boxexo}[title={\ding{45}\quad Exercice~14.6. TSANAGA \niveaumoyen}]
Que révèle l'extraction logique du téléphone de M.~NDONGO~?
\end{boxexo}

\begin{boxexo}[title={\ding{45}\quad Exercice~14.7. Géolocalisation \niveaumoyen}]
Que peut apporter la géolocalisation à l'enquête~?
\end{boxexo}

\begin{boxexo}[title={\ding{45}\quad Exercice~14.8. Tiers \niveaumoyen}]
Comment protéger les correspondants dont les messages figurent sur le
téléphone~?
\end{boxexo}

\begin{boxexo}[title={\ding{45}\quad Exercice~14.9. Procédure \niveauexpert}]
Rédigez la procédure de saisie et d'extraction mobile pour TSANAGA.
\end{boxexo}

\begin{boxexo}[title={\ding{45}\quad Exercice~14.10. Faisceau \niveauexpert}]
Intégrez les messages mobiles au faisceau d'indices de l'affaire.
\end{boxexo}

\bigskip

\begin{boxinfo}
\textbf{Pour aller plus loin.}
\begin{itemize}
    \item NIST SP~800-101r1~\cite{nist800101}~; ISO/IEC 27037~\cite{iso27037}.
\end{itemize}
\end{boxinfo}

\begin{boxinfo}
\textbf{Transition.} Les données vivent aussi hors de l'appareil~: le
Chapitre~\ref{chap:nuage} traite du forensique du nuage et des applications.
\end{boxinfo}
